% Options for packages loaded elsewhere
\PassOptionsToPackage{unicode}{hyperref}
\PassOptionsToPackage{hyphens}{url}
\documentclass[
]{article}
\usepackage{xcolor}
\usepackage[margin=1in]{geometry}
\usepackage{amsmath,amssymb}
\setcounter{secnumdepth}{-\maxdimen} % remove section numbering
\usepackage{iftex}
\ifPDFTeX
  \usepackage[T1]{fontenc}
  \usepackage[utf8]{inputenc}
  \usepackage{textcomp} % provide euro and other symbols
\else % if luatex or xetex
  \usepackage{unicode-math} % this also loads fontspec
  \defaultfontfeatures{Scale=MatchLowercase}
  \defaultfontfeatures[\rmfamily]{Ligatures=TeX,Scale=1}
\fi
\usepackage{lmodern}
\ifPDFTeX\else
  % xetex/luatex font selection
\fi
% Use upquote if available, for straight quotes in verbatim environments
\IfFileExists{upquote.sty}{\usepackage{upquote}}{}
\IfFileExists{microtype.sty}{% use microtype if available
  \usepackage[]{microtype}
  \UseMicrotypeSet[protrusion]{basicmath} % disable protrusion for tt fonts
}{}
\makeatletter
\@ifundefined{KOMAClassName}{% if non-KOMA class
  \IfFileExists{parskip.sty}{%
    \usepackage{parskip}
  }{% else
    \setlength{\parindent}{0pt}
    \setlength{\parskip}{6pt plus 2pt minus 1pt}}
}{% if KOMA class
  \KOMAoptions{parskip=half}}
\makeatother
\usepackage{color}
\usepackage{fancyvrb}
\newcommand{\VerbBar}{|}
\newcommand{\VERB}{\Verb[commandchars=\\\{\}]}
\DefineVerbatimEnvironment{Highlighting}{Verbatim}{commandchars=\\\{\}}
% Add ',fontsize=\small' for more characters per line
\usepackage{framed}
\definecolor{shadecolor}{RGB}{248,248,248}
\newenvironment{Shaded}{\begin{snugshade}}{\end{snugshade}}
\newcommand{\AlertTok}[1]{\textcolor[rgb]{0.94,0.16,0.16}{#1}}
\newcommand{\AnnotationTok}[1]{\textcolor[rgb]{0.56,0.35,0.01}{\textbf{\textit{#1}}}}
\newcommand{\AttributeTok}[1]{\textcolor[rgb]{0.13,0.29,0.53}{#1}}
\newcommand{\BaseNTok}[1]{\textcolor[rgb]{0.00,0.00,0.81}{#1}}
\newcommand{\BuiltInTok}[1]{#1}
\newcommand{\CharTok}[1]{\textcolor[rgb]{0.31,0.60,0.02}{#1}}
\newcommand{\CommentTok}[1]{\textcolor[rgb]{0.56,0.35,0.01}{\textit{#1}}}
\newcommand{\CommentVarTok}[1]{\textcolor[rgb]{0.56,0.35,0.01}{\textbf{\textit{#1}}}}
\newcommand{\ConstantTok}[1]{\textcolor[rgb]{0.56,0.35,0.01}{#1}}
\newcommand{\ControlFlowTok}[1]{\textcolor[rgb]{0.13,0.29,0.53}{\textbf{#1}}}
\newcommand{\DataTypeTok}[1]{\textcolor[rgb]{0.13,0.29,0.53}{#1}}
\newcommand{\DecValTok}[1]{\textcolor[rgb]{0.00,0.00,0.81}{#1}}
\newcommand{\DocumentationTok}[1]{\textcolor[rgb]{0.56,0.35,0.01}{\textbf{\textit{#1}}}}
\newcommand{\ErrorTok}[1]{\textcolor[rgb]{0.64,0.00,0.00}{\textbf{#1}}}
\newcommand{\ExtensionTok}[1]{#1}
\newcommand{\FloatTok}[1]{\textcolor[rgb]{0.00,0.00,0.81}{#1}}
\newcommand{\FunctionTok}[1]{\textcolor[rgb]{0.13,0.29,0.53}{\textbf{#1}}}
\newcommand{\ImportTok}[1]{#1}
\newcommand{\InformationTok}[1]{\textcolor[rgb]{0.56,0.35,0.01}{\textbf{\textit{#1}}}}
\newcommand{\KeywordTok}[1]{\textcolor[rgb]{0.13,0.29,0.53}{\textbf{#1}}}
\newcommand{\NormalTok}[1]{#1}
\newcommand{\OperatorTok}[1]{\textcolor[rgb]{0.81,0.36,0.00}{\textbf{#1}}}
\newcommand{\OtherTok}[1]{\textcolor[rgb]{0.56,0.35,0.01}{#1}}
\newcommand{\PreprocessorTok}[1]{\textcolor[rgb]{0.56,0.35,0.01}{\textit{#1}}}
\newcommand{\RegionMarkerTok}[1]{#1}
\newcommand{\SpecialCharTok}[1]{\textcolor[rgb]{0.81,0.36,0.00}{\textbf{#1}}}
\newcommand{\SpecialStringTok}[1]{\textcolor[rgb]{0.31,0.60,0.02}{#1}}
\newcommand{\StringTok}[1]{\textcolor[rgb]{0.31,0.60,0.02}{#1}}
\newcommand{\VariableTok}[1]{\textcolor[rgb]{0.00,0.00,0.00}{#1}}
\newcommand{\VerbatimStringTok}[1]{\textcolor[rgb]{0.31,0.60,0.02}{#1}}
\newcommand{\WarningTok}[1]{\textcolor[rgb]{0.56,0.35,0.01}{\textbf{\textit{#1}}}}
\usepackage{graphicx}
\makeatletter
\newsavebox\pandoc@box
\newcommand*\pandocbounded[1]{% scales image to fit in text height/width
  \sbox\pandoc@box{#1}%
  \Gscale@div\@tempa{\textheight}{\dimexpr\ht\pandoc@box+\dp\pandoc@box\relax}%
  \Gscale@div\@tempb{\linewidth}{\wd\pandoc@box}%
  \ifdim\@tempb\p@<\@tempa\p@\let\@tempa\@tempb\fi% select the smaller of both
  \ifdim\@tempa\p@<\p@\scalebox{\@tempa}{\usebox\pandoc@box}%
  \else\usebox{\pandoc@box}%
  \fi%
}
% Set default figure placement to htbp
\def\fps@figure{htbp}
\makeatother
\setlength{\emergencystretch}{3em} % prevent overfull lines
\providecommand{\tightlist}{%
  \setlength{\itemsep}{0pt}\setlength{\parskip}{0pt}}
\usepackage{bookmark}
\IfFileExists{xurl.sty}{\usepackage{xurl}}{} % add URL line breaks if available
\urlstyle{same}
\hypersetup{
  pdftitle={Term Project},
  pdfauthor={Liangmi Zhang},
  hidelinks,
  pdfcreator={LaTeX via pandoc}}

\title{Term Project}
\author{Liangmi Zhang}
\date{November 24, 2025}

\begin{document}
\maketitle

\section{An Analysis of Coronary Heart
Disease}\label{an-analysis-of-coronary-heart-disease}

\subsection{Header 2}\label{header-2}

\subsection{Header 2}\label{header-2-1}

\subsection{Logistic Regression}\label{logistic-regression}

In this section, we try to predict if a person has a coronary heart
disease.

\subsubsection{Logistic Regression Without Using
Weights}\label{logistic-regression-without-using-weights}

I ignored the various weights in the dataset because I couldn't figure
out what they meant and how to use them. However, my models produce AUC
around 5. I thought the poor result was because of the missing weights
so I decided to use weights to improve the result. However, when I were
preparing the RMD file, I found out that the low AUC was because of a
bug - in my original code, I intended to train data on the training set
that was not defined in the script, and the script was using another
dataframe called \textbf{train} from the global environment. After
correcting the bug, the ROC looks normal. The corrected result is
provided below.

Training with 3 predictors

\begin{Shaded}
\begin{Highlighting}[]
\CommentTok{\# {-}{-}{-}{-}{-}model with multiple predictors{-}{-}{-}{-}{-}{-}{-}{-}{-}{-}{-}.}
\CommentTok{\# 1. Fit on training set}
\NormalTok{model }\OtherTok{\textless{}{-}} \FunctionTok{glm}\NormalTok{(}
\NormalTok{  coronary }\SpecialCharTok{\textasciitilde{}}\NormalTok{ high\_density\_level }\SpecialCharTok{+}\NormalTok{ low\_density\_level }\SpecialCharTok{+}\NormalTok{ triglyceride\_level,}
  \AttributeTok{family =}\NormalTok{ binomial,}
  \AttributeTok{data =}\NormalTok{ train}
\NormalTok{)}

\CommentTok{\# 2. Predict on *validation* set}
\NormalTok{validation}\SpecialCharTok{$}\NormalTok{prob }\OtherTok{\textless{}{-}} \FunctionTok{predict}\NormalTok{(}
\NormalTok{  model,}
  \AttributeTok{newdata =}\NormalTok{ validation,}
  \AttributeTok{type =} \StringTok{"response"}
\NormalTok{)}

\CommentTok{\# 3. ROC on validation set}

\NormalTok{g }\OtherTok{\textless{}{-}} \FunctionTok{roc}\NormalTok{(}\AttributeTok{response =}\NormalTok{ validation}\SpecialCharTok{$}\NormalTok{coronary,}
         \AttributeTok{predictor =}\NormalTok{ validation}\SpecialCharTok{$}\NormalTok{prob)}
\end{Highlighting}
\end{Shaded}

\begin{verbatim}
## Setting levels: control = 0, case = 1
\end{verbatim}

\begin{verbatim}
## Setting direction: controls < cases
\end{verbatim}

\begin{Shaded}
\begin{Highlighting}[]
\NormalTok{g}
\end{Highlighting}
\end{Shaded}

\begin{verbatim}
## 
## Call:
## roc.default(response = validation$coronary, predictor = validation$prob)
## 
## Data: validation$prob in 961 controls (validation$coronary 0) < 60 cases (validation$coronary 1).
## Area under the curve: 0.7247
\end{verbatim}

\begin{Shaded}
\begin{Highlighting}[]
\FunctionTok{plot}\NormalTok{(g)}
\end{Highlighting}
\end{Shaded}

\pandocbounded{\includegraphics[keepaspectratio]{termProject_files/figure-latex/fittingModelWithoutWeightMultiplePredictors-1.pdf}}

Training with one predictor.

\begin{Shaded}
\begin{Highlighting}[]
\CommentTok{\# {-}{-}{-}{-}{-}{-}{-}{-}{-} 1. Fit model on TRAIN only {-}{-}{-}{-}{-}{-}{-}{-}{-}{-}{-}{-}{-}{-}}

\NormalTok{model }\OtherTok{\textless{}{-}} \FunctionTok{glm}\NormalTok{(}
\NormalTok{  coronary }\SpecialCharTok{\textasciitilde{}}\NormalTok{ low\_density\_level,}
  \AttributeTok{family =}\NormalTok{ binomial,}
  \AttributeTok{data =}\NormalTok{ train}
\NormalTok{)}

\CommentTok{\# {-}{-}{-}{-}{-}{-}{-}{-}{-} 2. Predict on VALIDATION only {-}{-}{-}{-}{-}{-}{-}{-}{-}{-}{-}}

\NormalTok{validation}\SpecialCharTok{$}\NormalTok{prob }\OtherTok{\textless{}{-}} \FunctionTok{predict}\NormalTok{(}
\NormalTok{  model,}
  \AttributeTok{newdata =}\NormalTok{ validation,}
  \AttributeTok{type =} \StringTok{"response"}
\NormalTok{)}

\CommentTok{\# {-}{-}{-}{-}{-}{-}{-}{-}{-} 3. Compute ROC on VALIDATION {-}{-}{-}{-}{-}{-}{-}{-}{-}{-}{-}}

\FunctionTok{library}\NormalTok{(pROC)}

\NormalTok{g }\OtherTok{\textless{}{-}} \FunctionTok{roc}\NormalTok{(}
  \AttributeTok{response  =}\NormalTok{ validation}\SpecialCharTok{$}\NormalTok{coronary,}
  \AttributeTok{predictor =}\NormalTok{ validation}\SpecialCharTok{$}\NormalTok{prob}
\NormalTok{)}
\end{Highlighting}
\end{Shaded}

\begin{verbatim}
## Setting levels: control = 0, case = 1
\end{verbatim}

\begin{verbatim}
## Setting direction: controls < cases
\end{verbatim}

\begin{Shaded}
\begin{Highlighting}[]
\NormalTok{g}
\end{Highlighting}
\end{Shaded}

\begin{verbatim}
## 
## Call:
## roc.default(response = validation$coronary, predictor = validation$prob)
## 
## Data: validation$prob in 961 controls (validation$coronary 0) < 60 cases (validation$coronary 1).
## Area under the curve: 0.7416
\end{verbatim}

\begin{Shaded}
\begin{Highlighting}[]
\FunctionTok{plot}\NormalTok{(g)}
\end{Highlighting}
\end{Shaded}

\pandocbounded{\includegraphics[keepaspectratio]{termProject_files/figure-latex/fittingModelWithoutWeightOnePredictor-1.pdf}}

\subsubsection{Logistic Regression With
Weights}\label{logistic-regression-with-weights}

Because of low AUC value in the initial trials, I decided to incooporate
weights in the regression. In order to learn how to use them, I studied
NAHNAES's sample design, a concept that I haven't studied before and
struggled to understand. Another challenge is how to choose a weight
since there are many weights in the dataset. It took me a while to
understand the criteria - choose the weight for the smallest subsample
if multiple subsamples are used in the analysis. In my case, I needed to
decide which subsample is smaller because I have two: the phlebotomy
subsample(people who gave blood specimen) and the fasting subsample
(people who are asked to fast). It took me several hours just to decide
on this. Finally, I realized that fasting subsample was smaller because
people in the sample fasted first and then gave blood specimen.

\paragraph{Logistic Regression With Using Weight\_2(Fasting Subsample
Weight)}\label{logistic-regression-with-using-weight_2fasting-subsample-weight}

In the beginning,the models couldn't converge, no matter the number of
predictors that I used. This turned out to be another bug that I didn't
know. I thought the failure was due to lack of datapoints in the
subsample because the fasting subsample is the smallest subsample.
However, somehow, the models could converge after some changes that I
made, but I forgot to track where I changed. The models are presented
below.

Trying to fit a model with 3 predictors. This model won't converge.

\begin{Shaded}
\begin{Highlighting}[]
\CommentTok{\# {-}{-}{-}{-}{-}{-}{-}{-}{-} model with 3 predictors {-}{-}{-}{-}{-}{-}{-}{-}{-}{-}{-}{-}{-}{-}{-}{-}{-}{-}{-}{-}{-}{-}{-}}
\CommentTok{\# train model on the train set}
\NormalTok{model1 }\OtherTok{\textless{}{-}} \FunctionTok{svyglm}\NormalTok{(}
\NormalTok{  coronary }\SpecialCharTok{\textasciitilde{}}\NormalTok{ high\_density\_level }\SpecialCharTok{+}\NormalTok{low\_density\_level }\SpecialCharTok{+}\NormalTok{triglyceride\_level  ,}
  \AttributeTok{design =}\NormalTok{ design\_train\_sub,}
  \AttributeTok{family =} \FunctionTok{quasibinomial}\NormalTok{()}
\NormalTok{)}

\CommentTok{\# predict on the validation set.}

\NormalTok{val\_df }\OtherTok{\textless{}{-}}\NormalTok{ design\_validation\_sub}\SpecialCharTok{$}\NormalTok{variables}
\NormalTok{val\_df}\SpecialCharTok{$}\NormalTok{predicted }\OtherTok{\textless{}{-}} \FunctionTok{predict}\NormalTok{(model1, }\AttributeTok{newdata =}\NormalTok{ val\_df, }\AttributeTok{type =} \StringTok{"response"}\NormalTok{)}

\NormalTok{roc\_val\_three }\OtherTok{\textless{}{-}} \FunctionTok{roc}\NormalTok{(}
  \AttributeTok{response =}\NormalTok{ val\_df}\SpecialCharTok{$}\NormalTok{coronary,}
  \AttributeTok{predictor =}\NormalTok{ val\_df}\SpecialCharTok{$}\NormalTok{predicted,}
  \AttributeTok{weights =}\NormalTok{ val\_df}\SpecialCharTok{$}\NormalTok{weight\_2}
\NormalTok{)}
\end{Highlighting}
\end{Shaded}

\begin{verbatim}
## Setting levels: control = 0, case = 1
\end{verbatim}

\begin{verbatim}
## Setting direction: controls < cases
\end{verbatim}

\begin{Shaded}
\begin{Highlighting}[]
\NormalTok{roc\_val\_three}
\end{Highlighting}
\end{Shaded}

\begin{verbatim}
## 
## Call:
## roc.default(response = val_df$coronary, predictor = val_df$predicted,     weights = val_df$weight_2)
## 
## Data: val_df$predicted in 970 controls (val_df$coronary 0) < 59 cases (val_df$coronary 1).
## Area under the curve: 0.7521
\end{verbatim}

\begin{Shaded}
\begin{Highlighting}[]
\FunctionTok{plot}\NormalTok{(roc\_val\_three)}
\end{Highlighting}
\end{Shaded}

\pandocbounded{\includegraphics[keepaspectratio]{termProject_files/figure-latex/weight2with3predictors-1.pdf}}

Fit another model with low\_density\_cholestrol.

\begin{Shaded}
\begin{Highlighting}[]
\CommentTok{\# {-}{-}{-}{-}{-}{-}{-}{-}{-} model one predictor : low density level {-}{-}{-}{-}{-}{-}{-}{-}{-}{-}{-}{-}{-}{-}{-}{-}{-}{-}{-}{-}{-}{-}{-}}

\CommentTok{\# train model on the train set}
\NormalTok{model5 }\OtherTok{\textless{}{-}} \FunctionTok{svyglm}\NormalTok{(}
\NormalTok{  coronary }\SpecialCharTok{\textasciitilde{}}\NormalTok{ low\_density\_level ,}
  \AttributeTok{design =}\NormalTok{ design\_train\_sub,}
  \AttributeTok{family =} \FunctionTok{quasibinomial}\NormalTok{()}
\NormalTok{)}

\FunctionTok{summary}\NormalTok{(model5)}
\end{Highlighting}
\end{Shaded}

\begin{verbatim}
## 
## Call:
## svyglm(formula = coronary ~ low_density_level, design = design_train_sub, 
##     family = quasibinomial())
## 
## Survey design:
## subset(design_train, !is.na(coronary) & !is.na(high_density_level) & 
##     !is.na(low_density_level) & !is.na(triglyceride_level))
## 
## Coefficients:
##                    Estimate Std. Error t value Pr(>|t|)    
## (Intercept)        0.011181   0.541952   0.021 0.983832    
## low_density_level -0.033445   0.007259  -4.607 0.000407 ***
## ---
## Signif. codes:  0 '***' 0.001 '**' 0.01 '*' 0.05 '.' 0.1 ' ' 1
## 
## (Dispersion parameter for quasibinomial family taken to be 1.422759)
## 
## Number of Fisher Scoring iterations: 7
\end{verbatim}

\begin{Shaded}
\begin{Highlighting}[]
\CommentTok{\# predict on the validation set.}

\NormalTok{val\_df }\OtherTok{\textless{}{-}}\NormalTok{ design\_validation\_sub}\SpecialCharTok{$}\NormalTok{variables}
\NormalTok{val\_df}\SpecialCharTok{$}\NormalTok{predicted }\OtherTok{\textless{}{-}} \FunctionTok{predict}\NormalTok{(model5, }\AttributeTok{newdata =}\NormalTok{ val\_df, }\AttributeTok{type =} \StringTok{"response"}\NormalTok{)}

\NormalTok{roc\_val\_LDL }\OtherTok{\textless{}{-}} \FunctionTok{roc}\NormalTok{(}
  \AttributeTok{response =}\NormalTok{ val\_df}\SpecialCharTok{$}\NormalTok{coronary,}
  \AttributeTok{predictor =}\NormalTok{ val\_df}\SpecialCharTok{$}\NormalTok{predicted,}
  \AttributeTok{weights =}\NormalTok{ val\_df}\SpecialCharTok{$}\NormalTok{weight\_2}
\NormalTok{)}
\end{Highlighting}
\end{Shaded}

\begin{verbatim}
## Setting levels: control = 0, case = 1
\end{verbatim}

\begin{verbatim}
## Setting direction: controls < cases
\end{verbatim}

\begin{Shaded}
\begin{Highlighting}[]
\NormalTok{roc\_val\_LDL}
\end{Highlighting}
\end{Shaded}

\begin{verbatim}
## 
## Call:
## roc.default(response = val_df$coronary, predictor = val_df$predicted,     weights = val_df$weight_2)
## 
## Data: val_df$predicted in 970 controls (val_df$coronary 0) < 59 cases (val_df$coronary 1).
## Area under the curve: 0.7417
\end{verbatim}

\begin{Shaded}
\begin{Highlighting}[]
\FunctionTok{plot}\NormalTok{(roc\_val\_LDL)}
\end{Highlighting}
\end{Shaded}

\pandocbounded{\includegraphics[keepaspectratio]{termProject_files/figure-latex/weight2with1predictor-1.pdf}}

\paragraph{Logistic Regression With Using Weight\_1(Phlebotomy Subsample
Weight)}\label{logistic-regression-with-using-weight_1phlebotomy-subsample-weight}

Since the previous method fails, the Phlebotomy weight is used since it
is the weight for the second smallest subsample.

\subparagraph{Logistic Regression With Using Weight\_1 and Two
Predictors}\label{logistic-regression-with-using-weight_1-and-two-predictors}

We fit a model using high\_density\_level and low\_density\_level.

\begin{Shaded}
\begin{Highlighting}[]
\CommentTok{\# {-}{-}{-}{-}{-}{-}{-}{-}{-} model with two predictors {-}{-}{-}{-}{-}{-}{-}{-}{-}{-}{-}{-}{-}{-}{-}{-}{-}{-}{-}{-}{-}{-}{-}}
\CommentTok{\# train model on the train set}
\NormalTok{model }\OtherTok{\textless{}{-}} \FunctionTok{svyglm}\NormalTok{(}
\NormalTok{  coronary }\SpecialCharTok{\textasciitilde{}}\NormalTok{ high\_density\_level }\SpecialCharTok{+}\NormalTok{low\_density\_level ,}
  \AttributeTok{design =}\NormalTok{ design\_train\_sub,}
  \AttributeTok{family =} \FunctionTok{quasibinomial}\NormalTok{()}
\NormalTok{)}

\FunctionTok{summary}\NormalTok{(model)}
\end{Highlighting}
\end{Shaded}

\begin{verbatim}
## 
## Call:
## svyglm(formula = coronary ~ high_density_level + low_density_level, 
##     design = design_train_sub, family = quasibinomial())
## 
## Survey design:
## subset(design_train, !is.na(coronary) & !is.na(high_density_level) & 
##     !is.na(low_density_level))
## 
## Coefficients:
##                     Estimate Std. Error t value Pr(>|t|)    
## (Intercept)         1.338345   0.664742   2.013 0.065262 .  
## high_density_level -0.023914   0.009549  -2.504 0.026367 *  
## low_density_level  -0.031691   0.007010  -4.521 0.000575 ***
## ---
## Signif. codes:  0 '***' 0.001 '**' 0.01 '*' 0.05 '.' 0.1 ' ' 1
## 
## (Dispersion parameter for quasibinomial family taken to be 1.38474)
## 
## Number of Fisher Scoring iterations: 6
\end{verbatim}

\begin{Shaded}
\begin{Highlighting}[]
\CommentTok{\# predict on the validation set.}

\NormalTok{val\_df }\OtherTok{\textless{}{-}}\NormalTok{ design\_validation\_sub}\SpecialCharTok{$}\NormalTok{variables}
\NormalTok{val\_df}\SpecialCharTok{$}\NormalTok{predicted }\OtherTok{\textless{}{-}} \FunctionTok{predict}\NormalTok{(model, }\AttributeTok{newdata =}\NormalTok{ val\_df, }\AttributeTok{type =} \StringTok{"response"}\NormalTok{)}

\NormalTok{roc\_val\_both }\OtherTok{\textless{}{-}} \FunctionTok{roc}\NormalTok{(}
  \AttributeTok{response =}\NormalTok{ val\_df}\SpecialCharTok{$}\NormalTok{coronary,}
  \AttributeTok{predictor =}\NormalTok{ val\_df}\SpecialCharTok{$}\NormalTok{predicted,}
  \AttributeTok{weights =}\NormalTok{ val\_df}\SpecialCharTok{$}\NormalTok{weight\_1}
\NormalTok{)}
\end{Highlighting}
\end{Shaded}

\begin{verbatim}
## Setting levels: control = 0, case = 1
\end{verbatim}

\begin{verbatim}
## Setting direction: controls < cases
\end{verbatim}

\begin{Shaded}
\begin{Highlighting}[]
\NormalTok{roc\_val\_both}
\end{Highlighting}
\end{Shaded}

\begin{verbatim}
## 
## Call:
## roc.default(response = val_df$coronary, predictor = val_df$predicted,     weights = val_df$weight_1)
## 
## Data: val_df$predicted in 941 controls (val_df$coronary 0) < 59 cases (val_df$coronary 1).
## Area under the curve: 0.7381
\end{verbatim}

\begin{Shaded}
\begin{Highlighting}[]
\FunctionTok{plot}\NormalTok{(roc\_val\_both)}
\end{Highlighting}
\end{Shaded}

\pandocbounded{\includegraphics[keepaspectratio]{termProject_files/figure-latex/regressionWithWeight_1With2Predictors-1.pdf}}

\subparagraph{Logistic Regression With Using Weight\_1 and One
Predictor}\label{logistic-regression-with-using-weight_1-and-one-predictor}

With high\_density\_level

We fit another model with high\_density\_level.

\begin{Shaded}
\begin{Highlighting}[]
\CommentTok{\# train model on the train set}
\NormalTok{model }\OtherTok{\textless{}{-}} \FunctionTok{svyglm}\NormalTok{(}
\NormalTok{  coronary }\SpecialCharTok{\textasciitilde{}}\NormalTok{ high\_density\_level ,}
  \AttributeTok{design =}\NormalTok{ design\_train\_sub,}
  \AttributeTok{family =} \FunctionTok{quasibinomial}\NormalTok{()}
\NormalTok{)}

\FunctionTok{summary}\NormalTok{(model)}
\end{Highlighting}
\end{Shaded}

\begin{verbatim}
## 
## Call:
## svyglm(formula = coronary ~ high_density_level, design = design_train_sub, 
##     family = quasibinomial())
## 
## Survey design:
## subset(design_train, !is.na(coronary) & !is.na(high_density_level) & 
##     !is.na(low_density_level))
## 
## Coefficients:
##                     Estimate Std. Error t value Pr(>|t|)   
## (Intercept)        -1.239641   0.468932  -2.644  0.01927 * 
## high_density_level -0.032451   0.008964  -3.620  0.00278 **
## ---
## Signif. codes:  0 '***' 0.001 '**' 0.01 '*' 0.05 '.' 0.1 ' ' 1
## 
## (Dispersion parameter for quasibinomial family taken to be 0.9851106)
## 
## Number of Fisher Scoring iterations: 6
\end{verbatim}

\begin{Shaded}
\begin{Highlighting}[]
\CommentTok{\# predict on the validation set.}

\NormalTok{val\_df }\OtherTok{\textless{}{-}}\NormalTok{ design\_validation\_sub}\SpecialCharTok{$}\NormalTok{variables}
\NormalTok{val\_df}\SpecialCharTok{$}\NormalTok{predicted }\OtherTok{\textless{}{-}} \FunctionTok{predict}\NormalTok{(model, }\AttributeTok{newdata =}\NormalTok{ val\_df, }\AttributeTok{type =} \StringTok{"response"}\NormalTok{)}

\NormalTok{roc\_val\_HDL }\OtherTok{\textless{}{-}} \FunctionTok{roc}\NormalTok{(}
  \AttributeTok{response =}\NormalTok{ val\_df}\SpecialCharTok{$}\NormalTok{coronary,}
  \AttributeTok{predictor =}\NormalTok{ val\_df}\SpecialCharTok{$}\NormalTok{predicted,}
  \AttributeTok{weights =}\NormalTok{ val\_df}\SpecialCharTok{$}\NormalTok{weight\_1}
\NormalTok{)}
\end{Highlighting}
\end{Shaded}

\begin{verbatim}
## Setting levels: control = 0, case = 1
\end{verbatim}

\begin{verbatim}
## Setting direction: controls < cases
\end{verbatim}

\begin{Shaded}
\begin{Highlighting}[]
\NormalTok{roc\_val\_HDL}
\end{Highlighting}
\end{Shaded}

\begin{verbatim}
## 
## Call:
## roc.default(response = val_df$coronary, predictor = val_df$predicted,     weights = val_df$weight_1)
## 
## Data: val_df$predicted in 941 controls (val_df$coronary 0) < 59 cases (val_df$coronary 1).
## Area under the curve: 0.5673
\end{verbatim}

\begin{Shaded}
\begin{Highlighting}[]
\FunctionTok{plot}\NormalTok{(roc\_val\_HDL)}
\end{Highlighting}
\end{Shaded}

\pandocbounded{\includegraphics[keepaspectratio]{termProject_files/figure-latex/regressionWithWeight_1With1HDL-1.pdf}}

With low\_density\_level

We fit another model with low\_density\_level.

\begin{Shaded}
\begin{Highlighting}[]
\CommentTok{\# train model on the train set}
\NormalTok{model }\OtherTok{\textless{}{-}} \FunctionTok{svyglm}\NormalTok{(}
\NormalTok{  coronary }\SpecialCharTok{\textasciitilde{}}\NormalTok{ low\_density\_level ,}
  \AttributeTok{design =}\NormalTok{ design\_train\_sub,}
  \AttributeTok{family =} \FunctionTok{quasibinomial}\NormalTok{()}
\NormalTok{)}

\FunctionTok{summary}\NormalTok{(model)}
\end{Highlighting}
\end{Shaded}

\begin{verbatim}
## 
## Call:
## svyglm(formula = coronary ~ low_density_level, design = design_train_sub, 
##     family = quasibinomial())
## 
## Survey design:
## subset(design_train, !is.na(coronary) & !is.na(high_density_level) & 
##     !is.na(low_density_level))
## 
## Coefficients:
##                    Estimate Std. Error t value Pr(>|t|)    
## (Intercept)        0.245026   0.605324   0.405 0.691752    
## low_density_level -0.033268   0.007198  -4.622 0.000396 ***
## ---
## Signif. codes:  0 '***' 0.001 '**' 0.01 '*' 0.05 '.' 0.1 ' ' 1
## 
## (Dispersion parameter for quasibinomial family taken to be 1.518632)
## 
## Number of Fisher Scoring iterations: 6
\end{verbatim}

\begin{Shaded}
\begin{Highlighting}[]
\CommentTok{\# predict on the validation set.}

\NormalTok{val\_df }\OtherTok{\textless{}{-}}\NormalTok{ design\_validation\_sub}\SpecialCharTok{$}\NormalTok{variables}
\NormalTok{val\_df}\SpecialCharTok{$}\NormalTok{predicted }\OtherTok{\textless{}{-}} \FunctionTok{predict}\NormalTok{(model, }\AttributeTok{newdata =}\NormalTok{ val\_df, }\AttributeTok{type =} \StringTok{"response"}\NormalTok{)}

\NormalTok{roc\_val\_LDL }\OtherTok{\textless{}{-}} \FunctionTok{roc}\NormalTok{(}
  \AttributeTok{response =}\NormalTok{ val\_df}\SpecialCharTok{$}\NormalTok{coronary,}
  \AttributeTok{predictor =}\NormalTok{ val\_df}\SpecialCharTok{$}\NormalTok{predicted,}
  \AttributeTok{weights =}\NormalTok{ val\_df}\SpecialCharTok{$}\NormalTok{weight\_1}
\NormalTok{)}
\end{Highlighting}
\end{Shaded}

\begin{verbatim}
## Setting levels: control = 0, case = 1
\end{verbatim}

\begin{verbatim}
## Setting direction: controls < cases
\end{verbatim}

\begin{Shaded}
\begin{Highlighting}[]
\NormalTok{roc\_val\_LDL}
\end{Highlighting}
\end{Shaded}

\begin{verbatim}
## 
## Call:
## roc.default(response = val_df$coronary, predictor = val_df$predicted,     weights = val_df$weight_1)
## 
## Data: val_df$predicted in 941 controls (val_df$coronary 0) < 59 cases (val_df$coronary 1).
## Area under the curve: 0.7357
\end{verbatim}

\begin{Shaded}
\begin{Highlighting}[]
\FunctionTok{plot}\NormalTok{(roc\_val\_LDL)}
\end{Highlighting}
\end{Shaded}

\pandocbounded{\includegraphics[keepaspectratio]{termProject_files/figure-latex/regressionWithWeight_1With1LDL-1.pdf}}

From this model, we choose a threshhold that yields the highest
accuracy.

\begin{Shaded}
\begin{Highlighting}[]
\CommentTok{\# 1.choose the best threshold for this model and calculate the accuracy. }

\NormalTok{roc\_obj }\OtherTok{\textless{}{-}} \FunctionTok{roc}\NormalTok{(}\AttributeTok{response =}\NormalTok{ val\_df}\SpecialCharTok{$}\NormalTok{coronary, }\AttributeTok{predictor =}\NormalTok{ val\_df}\SpecialCharTok{$}\NormalTok{predicted)}
\end{Highlighting}
\end{Shaded}

\begin{verbatim}
## Setting levels: control = 0, case = 1
\end{verbatim}

\begin{verbatim}
## Setting direction: controls < cases
\end{verbatim}

\begin{Shaded}
\begin{Highlighting}[]
\NormalTok{best }\OtherTok{\textless{}{-}} \FunctionTok{coords}\NormalTok{(roc\_obj, }\StringTok{"best"}\NormalTok{, }\AttributeTok{best.method=}\StringTok{"closest.topleft"}\NormalTok{)}
\NormalTok{best}
\end{Highlighting}
\end{Shaded}

\begin{verbatim}
##    threshold specificity sensitivity
## 1 0.06707157   0.7651435   0.6271186
\end{verbatim}

\begin{Shaded}
\begin{Highlighting}[]
\NormalTok{best\_threshold }\OtherTok{\textless{}{-}} \FunctionTok{as.numeric}\NormalTok{(best}\SpecialCharTok{$}\NormalTok{threshold)}

\CommentTok{\# 2. predicted classes using optimal threshold}
\NormalTok{val\_df}\SpecialCharTok{$}\NormalTok{predicted\_labels }\OtherTok{\textless{}{-}} \FunctionTok{ifelse}\NormalTok{(val\_df}\SpecialCharTok{$}\NormalTok{predicted }\SpecialCharTok{\textgreater{}=}\NormalTok{ best\_threshold[], }\DecValTok{1}\NormalTok{, }\DecValTok{0}\NormalTok{)}

\CommentTok{\# 3. accuracy}
\NormalTok{accuracy }\OtherTok{\textless{}{-}} \FunctionTok{mean}\NormalTok{(val\_df}\SpecialCharTok{$}\NormalTok{predicted\_labels }\SpecialCharTok{==}\NormalTok{ val\_df}\SpecialCharTok{$}\NormalTok{coronary)}
\NormalTok{accuracy}
\end{Highlighting}
\end{Shaded}

\begin{verbatim}
## [1] 0.757
\end{verbatim}

\end{document}
